\definecolor{codegreen}{rgb}{0,0.6,0}
\definecolor{codegray}{rgb}{0.5,0.5,0.5}
\definecolor{codepurple}{rgb}{0.58,0,0.82}
\definecolor{backcolour}{rgb}{0.95,0.95,0.92}

\lstdefinestyle{mystyle}{
    backgroundcolor=\color{backcolour},
    commentstyle=\color{codegreen},
    keywordstyle=\color{magenta},
    numberstyle=\tiny\color{codegray},
    stringstyle=\color{codepurple},
    basicstyle=\ttfamily\footnotesize,
    breakatwhitespace=false,
    breaklines=true,
    captionpos=b,
    keepspaces=true,
    numbers=left,
    numbersep=5pt,
    showspaces=false,
    showstringspaces=false,
    showtabs=false,
    tabsize=2
}
\lstset{style=mystyle}


\section{Full Prompts}\label{appendix_full_prompt}
This section delineates the constituents of all system prompts utilized within AutoDAN-Turbo, providing a comprehensive explanation of the method parameters that formulate these prompts:

\setlength{\parindent}{20pt}\verb|goal|: This refers to the malicious behaviour we aim to address.

\setlength{\parindent}{20pt}\verb|strategies_list|: This is a list comprising of strategies retrieved through the 'Jailbreak Strategy Retrieve' process, as discussed in Section 3.3.

\setlength{\parindent}{20pt} \verb|prompt|: This is the jailbreak attack prompt P, produced by the attacker LLM.

\setlength{\parindent}{20pt} \verb|response|: This is the response R, generated by the target LLM.

\setlength{\parindent}{20pt} \verb|att1|: This is the jailbreak attack prompt $P_i$, produced by the attacker LLM in the i-th round of jailbreaking.

\setlength{\parindent}{20pt} \verb|res1|: This is the response $R_i$ from the target LLM during the i-th round of jailbreaking.

\setlength{\parindent}{20pt} \verb|att2|: This is the jailbreak attack prompt $P_{i+1}$, produced by the attacker LLM during the (i+1)-th round of jailbreaking.

\setlength{\parindent}{20pt} \verb|res2|: This is the response $R_{i+1}$ from the target LLM during the (i+1)-th round of jailbreaking.

\setlength{\parindent}{20pt} \verb|strategy_library|: This is the strategy library, housing all evolved jailbreak strategies to date.


\setlength{\parindent}{0pt}Employing this methodical approach, AutoDAN-Turbo ensures the delivery of clear, fluent, and precise malicious behaviour addressing strategies.
\subsection{Full Prompts for Configuring the Attacker LLM}\label{appendix_prompt_attacker}

\subsubsection{Initial Prompt Without Strategies Limitation}

In the process of constructing the Strategy Library, as outlined in section 3.2, this prompt serves as the system prompt for the attacker. Without any strategic limitations, the model is given maximum freedom to explore and unearth potentially successful jailbreaking strategies.
\lstinputlisting[language=Python]{sections/prompts/init_without_limitation.txt}


\subsubsection{Strategies Retrieval and Exploitation}
As highlighted in Section 3.3, we employ a specific system prompt to guide our attacker. This prompt enables the attacker to utilize the effective strategies retrieved (denoted as \verb|strategy_list|) to design a jailbreak attack prompt, specifically tailored towards a particular malicious behavior (denoted as \verb|goal|).
\lstinputlisting[language=Python]{sections/prompts/strategy_retrieval.txt}

\subsubsection{New Strategies Exploration}
When the strategies retrieved by our algorithm (referred to as the \verb|strategy_list|) are deemed ineffective, as shown in Table 1 of Section 3.1, we employ a unique approach to foster the evolution of new jailbreak strategies. We generate a system prompt for the attacker in the subsequent jailbreak attempt, designed to exclude the ineffective strategies and facilitate the emergence of more potent tactics. This approach ensures a dynamic evolution of jailbreak strategies, continuously refining our algorithm's effectiveness while maintaining a stringent exclusion of ineffective strategies.
\lstinputlisting[language=Python]{sections/prompts/find_new.txt}


\subsection{Full Prompts for Configuring the Scorer LLM}\label{appendix_prompt_scorer}
The scoring guidance for the Scorer LLM is embedded within the system prompt, which includes detailed scoring criteria. This enables the scorer LLM to accurately evaluate the effectiveness of the attacker prompt P and the corresponding target response R.
\lstinputlisting[language=Python]{sections/prompts/scorer.txt}


\subsection{Full Prompts for Configuring the Summarizer LLM}\label{appendix_prompt_summarizer}
As outlined in Section 3.2, during the process of building the strategy library, this prompt is utilized as the system prompt for the summarizer LLM, which guides it to discover and generate standardized jailbreak strategies in a logical, clear, and effective manner.
\lstinputlisting[language=Python]{sections/prompts/summarizer.txt}

